% =====================================================================
%  2.2 Likningssystemer — ekstra utfordring
% =====================================================================

\begin{infoboks}{Velg metode selv}
Arket bygger på 2.2 Likningssystemer. Vis regning og begrunnelse. Bruk eksakte svar.
Kontroller svaret ved å sette inn i alle de opprinnelige likningene.
\end{infoboks}

\nivaahode{2}

\oppg{U1}{2}{\uh}
Løs systemet med tre ukjente. Reduser først til to ukjente.
\[ \begin{cases} x + y + z = 6\\ x - y = 1\\ y + z = 4 \end{cases} \]

\oppg{U2}{2}{\mh}
På kino betaler 3 voksne og 2 barn 610 kr. 2 voksne og 4 barn betaler 620 kr.
\begin{tasks}(1)
  \task Sett opp et likningssystem og finn billettprisen for voksne og for barn.
  \task Tegn linjene i GeoGebra og kontroller skjæringspunktet.
\end{tasks}

\oppg{U3}{2}{\uh}
Løs systemet. Fjern brøkene først.
\[ \begin{cases} \dfrac{x}{2} + \dfrac{y}{3} = 4\\[2mm] x - y = 3 \end{cases} \]

\nivaahode{3}

\oppg{U4}{3}{\uh}
Se på systemet $\begin{cases} ax + by = c\\ dx + ey = f \end{cases}$ der $ae - bd \neq 0$.
\begin{tasks}(1)
  \task Vis at $x = \dfrac{ce - bf}{ae - bd}$. Gang første likning med $e$ og andre med $b$.
  \task Finn et tilsvarende uttrykk for $y$.
  \task Hva skjer når $ae - bd = 0$? Knytt svaret til stigningstall og linjer.
\end{tasks}

\oppg{U5}{3}{\prog}
Skriv en funksjon \lstinline[style=laereverk,numbers=none,frame=none]|los(a, b, c, d, e, f)| som bruker formlene fra U4.
Den skal returnere $(x, y)$. Når $ae - bd = 0$, skal den returnere teksten «ingen eller uendelig mange løsninger».
\begin{tasks}(1)
  \task Test funksjonen på $2x + 3y = 12$ og $5x - 2y = 11$, og på $2x - y = 3$ og $-4x + 2y = 6$.
  \task Et annet program prøver alle heltallspar med $0 \le x, y \le 20$. Nevn en fordel og en ulempe ved den strategien.
\end{tasks}

\oppg{U6}{3}{\uh}
Undersøk hvor mange skjæringspunkt en rett linje og parabelen $y = x^2$ kan ha.
\begin{tasks}(1)
  \task Løs $\begin{cases} y = x^2\\ y = 2x - 1 \end{cases}$. \hint{Flytt alt over på én side. Kjenner du igjen uttrykket?}
  \task Finn en linje som har to skjæringspunkt med parabelen, og en som ikke har noe. Skriv likningene og lag en skisse.
  \task Kan en linje ha tre skjæringspunkt med parabelen? Begrunn.
\end{tasks}

\begin{oppsummering}
Velg innsetting når en ukjent står alene, og addisjon når koeffisientene passer.
$0 = 0$ betyr samme linje, og en usann påstand betyr parallelle linjer.
Kontroller alltid svaret i alle likningene.
\end{oppsummering}
